\section{Introduction}
Self-evolving agents increasingly learn through external memory. A successful episode can become a reusable workflow; a failed episode can become a corrective reflection. Systems such as ReasoningBank convert trajectories into reasoning memories that influence future behavior \citep{ouyang2026reasoningbank}, while Agent Workflow Memory and Reflexion persist workflows or verbal feedback across trials \citep{wang2024awm,shinn2023reflexion}. A natural engineering shortcut is to treat the source outcome itself as a quality signal: trust memories written from successes and discount memories written from failures.

That shortcut confounds three objects that need not agree. First, \emph{memory content} is the actionable text presented to the agent. Second, \emph{process provenance} records how the source trajectory ended. Third, \emph{downstream marginal utility} asks whether reusing that memory helps a new task. A failed trajectory can yield a useful correction, while a successful trajectory can yield brittle or irrelevant guidance. Learning from Failure explicitly exploits failed computer-use trajectories \citep{sun2026learningfromfailure}, and multi-round skill-evolution results likewise show that useful updates can arise from failure feedback \citep{liu2026rethinkingfeedback}. Provenance therefore has no justified monotone sign a priori.

Existing memory research also makes it important to distinguish content from lineage. Content-centric systems such as A-MEM and MemRL improve memory organization, retrieval, or runtime utility \citep{xu2025amem,zhang2026memrl}. Other work explicitly represents execution provenance or shows that memory consolidation can obscure source lineage while preserving action triggers \citep{wang2026tracetotrust,xu2026laundering}. We do not claim that provenance-aware storage is new. Our narrower question is causal and decision-theoretic: \emph{conditional on the same actionable memory content and target context, does success/failure process provenance contain incremental information that improves downstream control?}

A released audit of self-evolving financial agents supplies a strong motivating association \citep{li2026audit}. In its ReasoningBank analysis, success-derived retrievals have reported accuracy 0.931 over 101 cases and failure-derived retrievals 0.647 over 34 cases; all ten reported persistent W$\rightarrow$W failures lie in the failure-derived stratum. Yet source-task difficulty can jointly affect acquisition outcome and later difficulty, so this L0 association does not identify a provenance effect. Nor does switching a success-conditioned versus failure-conditioned memory writer isolate provenance: the generated text, pragmatic register, and strategy emphasis may change together.

We therefore formalize a \emph{provenance-audit identification ladder}. L0 compares naturally occurring provenance strata. L1 changes an outcome-conditioned writer while fixing the downstream state; it identifies a writer-mode bundle, not provenance alone. L2 holds actionable memory content, retrieval membership, and order fixed and changes only executor-visible provenance metadata. L3 transports an identified intervention into the motivating source runtime. The ladder provides a constructive non-upgrade rule: an L1 effect cannot be promoted to L2 unless information identity is established.

Historical controlled evidence motivates but does not settle the cleaner estimand. On released WebArena evidence, an L1 terminal endpoint uses four independent task pairs and gives a success-writer minus failure-writer difference of $+0.1667$ with $p=0.0785$; four units cannot attain an exact one-sided $p<0.05$. A disjoint 23-task early-action endpoint gives a reference-action-match difference of $-0.0942$ with counterevidence-side $p=0.0792$. Because these are different estimands and the writer outputs differ in style and strategy emphasis, they support no provenance-only sign. Earlier L2 attempts correctly fail closed on support or runtime constraints and are not pooled into the experiment reported here.

Our main experiment prospectively executes L2 on a fresh MemRL/OSInteraction substrate. We first build an independent memory bank from the full 350-task frozen source split, yielding 176 success-derived and 174 failure-derived memories. Before any validation environment is reset, native MemRL retrieval is evaluated over fresh exact skill-signature clusters after excluding all historical primary/utilization clusters; 106/108 clusters have eligible support. Frozen hash ranking then selects 32 primary clusters and eight disjoint utilization clusters. The utilization first stage runs true, null, reversed, shuffled, and no-memory surfaces and passes its preregistered first-action gate, ruling out the trivial explanation that retrieved memory is behaviorally unused.

The primary test then compares two executor views on the same 32 clusters. Arm A exposes the frozen actionable content only. Arm B exposes the same retrieval membership, order, and actionable bytes plus truthful boolean source outcome. No retrieval is rerun between arms. With Qwen2.5-7B-Instruct, A succeeds on 15/32 tasks and B on 16/32, so the paired effect is $+0.03125$. There is only one discordant terminal pair; the exact two-sided sign test is $p=1.0$, the preregistered paired bootstrap interval is $[0,0.09375]$, and the prespecified $|\Delta|\geq0.15$ practical-relevance floor is not met. Importantly, provenance is not behaviorally invisible: A and B choose different first executable actions in 9/32 pairs and take different numbers of steps in 7/32, while terminal success differs in only 1/32. The pattern is therefore better described as \emph{behaviorally legible provenance with low terminal conversion} than as ``the agent ignores provenance.''

We prospectively replicate the exact same L2 surface with Meta-Llama-3.1-8B-Instruct while holding the source bank, frozen retrieval content/order, 32 primary clusters, renderer, evaluator, randomization, and analysis fixed. A model-specific parser qualification passes before validation exposure, and the Llama utilization first stage passes with memory-specific first-action divergence in 8/8 units versus 5/8 placebo divergences. All 64 primary arms then complete. The Llama paired terminal effect is exactly $0$: two clusters are B-only successes and two are A-only successes, giving four discordant pairs, exact $p=1.0$, and a paired bootstrap interval of $[-0.125,0.125]$. The 0.15 relevance floor is again not met. We report the Qwen and Llama estimates separately rather than pooling them. The cross-backbone replication therefore strengthens the bounded statement that raw provenance has low incremental terminal value on this exact-information MemRL/OSInteraction surface; it does not establish a universal null.

\paragraph{Contributions and closest-work boundary.}
We contribute (i) an L0--L3 audit contract that separates source association, writer-mode effects, exact-information provenance exposure, and source transport; (ii) a support-qualified, prospectively frozen L2 experiment with 350 source memories, an explicit utilization first stage, and 32 paired terminal units; (iii) an executor-backbone replication that changes only the model and reproduces the low-terminal-value regime on the same frozen information surface; and (iv) Provenance-Separated Memory Governance (PSMG), a design principle that retains lineage as auditable control-plane state but shrinks its decision influence when conditional information gain is unsupported. We do not claim a universal zero provenance effect, that failure-derived memories are generally harmful, that model-specific estimates may be pooled post hoc, or that PSMG efficacy has been demonstrated. Source-faithful L3 transport and a governor-versus-same-information-controller efficacy trial remain separate questions.
